% Actividad 1 contractualizada por AulaTeX: foro diagnóstico con audio.
\providecommand{\itescastudentname}{Martin Jonathan de la Cruz}
\providecommand{\itescastudentid}{Por confirmar}
\providecommand{\itescastudentemail}{Dato no proporcionado}
\providecommand{\itescafaculty}{Maestría en Gestión Administrativa}
\providecommand{\itescacoursename}{Fundamentos de Gestión Administrativa}
\providecommand{\itescacoursecode}{ACF--0901}
\providecommand{\itescadepartment}{Posgrado e Investigación}
\providecommand{\itescateachingfigure}{Docente de la materia por confirmar}
\providecommand{\itescaprogramlevel}{Maestría}
\providecommand{\itescaprogramperiod}{Agosto--noviembre de 2026}
\providecommand{\itescamodality}{En línea}
\providecommand{\itescalocation}{Cajeme, Sonora}
\providecommand{\itescaasset}{ITESCA/assets-itesca/itesca-monograma.png}
\providecommand{\itescaofficialasset}{ITESCA/assets-itesca/web/logo-itesca-contorno.png}
\providecommand{\itescasecondaryasset}{ITESCA/assets-itesca/itesca-monograma.png}
\providecommand{\itescabibliography}{ITESCA/maestria-en-gestion-administrativa/fundamentos-de-gestion-administrativa-mga/fundamentos-de-gestion-administrativa}

\def\itescadocumenttitle {La organización: identidad, límites y elementos de estudio}
\def\itescadocumentsubtitle {Actividad 1}
\def\itescadocumentsubject {Foro diagnóstico mediante audio}
\def\itescaactivityname {Actividad 1. Foro sobre la organización}
\def\itescasessionname {Unidad I}
\def\itescablockname {Fundamentos de Gestión Administrativa}
\def\itescamainproduct {Foro diagnóstico mediante audio}

\documentclass[spanish,letterpaper,oneside]{article}

\def\documenttitle {\itescadocumenttitle}
\def\documentsubtitle {\itescadocumentsubtitle}
\def\documentsubject {\itescadocumentsubject}
\def\documentauthor {\itescastudentname}
\def\coursename {\itescacoursename}
\def\coursecode {\itescacoursecode}
\def\universityname {Instituto Tecnológico Superior de Cajeme}
\def\universityfaculty {\itescafaculty}
\def\universitydepartment {\itescadepartment}
\def\universitydepartmentimage {\itescaofficialasset}
\def\universitydepartmentimagecfg {width=5.2cm}
\def\universitylocation {\itescalocation}

\def\authortable {
  \begin{tabular}{ll}
    \textbf{Alumno:} & \itescastudentname \\
    No. de control: & \itescastudentid \\
    Carrera: & \itescafaculty \\
    Semestre: & \itescaprogramperiod \\
    Asignatura: & \itescacoursename \\
    Docente: & \itescateachingfigure \\
    Modalidad: & \itescamodality \\
    Fecha de entrega: & \today \\
    Lugar: & \universitylocation
  \end{tabular}
}

\input{template}
\usepackage{helvet}
\renewcommand{\familydefault}{\sfdefault}
\usepackage[most]{tcolorbox}
\usepackage{attachfile}
\usepackage{enumitem}
\setcitestyle{numbers,open={[},close={]}}
\setlist[itemize]{leftmargin=*,nosep}

\definecolor{itescaRedDark}{HTML}{8F1117}
\definecolor{itescaRed}{HTML}{D70F18}
\definecolor{itescaCharcoal}{HTML}{141414}
\definecolor{itescaSilver}{HTML}{D9D9D9}
\definecolor{itescaPaper}{HTML}{F7F7F7}
\def\portraittitlecolor {itescaRedDark}

\attachfilesetup{color={0.843 0.059 0.094},print=false}
\newcommand{\foroCopyButton}[1]{%
  \textattachfile[appearance={},description={Participación del foro en texto copiable},mimetype=text/plain]{#1}{%
    \colorbox{itescaRed}{\small\textcolor{white}{\,Copiar participación\,}}%
  }%
}
\newtcolorbox{forobox}[1][]{%
  enhanced,breakable,
  colback=itescaPaper,colframe=itescaRed,
  boxrule=0pt,leftrule=3pt,arc=1.2mm,
  left=4mm,right=4mm,top=3mm,bottom=3mm,
  fonttitle=\bfseries\color{white},coltitle=white,colbacktitle=itescaRedDark,
  attach boxed title to top left={xshift=4mm,yshift=-2mm},
  boxed title style={colframe=itescaRedDark,arc=0.6mm},
  #1
}

\hypersetup{
  pdftitle={\itescadocumenttitle},
  pdfsubject={\itescadocumentsubject},
  pdfauthor={\itescastudentname},
  colorlinks=true,
  linkcolor=itescaRedDark,
  citecolor=itescaRedDark,
  urlcolor=itescaRedDark
}

\begin{document}
\templatePortrait
\templatePagecfg
\onehalfspacing

\begin{abstractd}
\textbf{Objetivo:} Explicar qué constituye una organización, delimitar qué agrupaciones no reúnen sus rasgos esenciales e identificar los elementos necesarios para desarrollar estudios organizacionales con enfoque sistémico y contextual.

\textbf{Producto:} Guion oral literal para una aportación de foro de máximo tres minutos y dos réplicas adaptables a las intervenciones de compañeros.

\textbf{Técnica didáctica:} Foro, técnica 42 del repertorio \textit{100 Técnicas Didácticas de Enseñanza y Aprendizaje}, orientada a aplicar información mediante diálogo crítico, reflexivo y respetuoso \citep{lopezForo2022}.
\end{abstractd}

\templateIndex
\templateFinalcfg

% Técnica didáctica: foro diagnóstico.
% Nivel cognitivo: comprender, distinguir, analizar y dialogar.
% Criterio de organización: concepto -> frontera -> sistema -> interacción.
\section{Introducción}

Distinguir una organización de cualquier agrupación humana es un problema administrativo relevante: si se confunde una coincidencia temporal con un sistema organizado, no pueden identificarse responsabilidades, procesos ni resultados. Una organización surge cuando personas y recursos se coordinan de manera relativamente estable alrededor de propósitos compartidos. Además, funciona como sistema abierto porque intercambia información, recursos y respuestas con su ambiente \citep{kastRosenzweig1972,kastRosenzweig1985}. Desde esa perspectiva, el análisis no debe aislar organigramas o individuos, sino explicar sus interdependencias. El producto central es un guion oral que delimita el concepto, señala lo que queda fuera de él y propone dimensiones para estudiar organizaciones.

\section{La organización como sistema humano y abierto}

\subsection{Preparación conceptual}

La estructura distribuye autoridad, tareas y coordinación, pero no agota la realidad organizacional. También intervienen personas, tecnología, cultura, poder, recursos, ambiente y mecanismos de retroalimentación. El diseño más pertinente depende del contexto y de las demandas que enfrenta la entidad \citep{bright2019,openstaxStructures2019}. Por ello, ni el carácter lucrativo ni la constitución jurídica son requisitos universales: una asociación comunitaria puede ser organización si sostiene objetivos, roles, coordinación y continuidad.

\subsection{Aportación preparada para el foro}

\begin{forobox}[title={Intervención en el foro sobre la organización}]
\hfill\foroCopyButton{ITESCA/maestria-en-gestion-administrativa/fundamentos-de-gestion-administrativa-mga/foro-participacion-Actividad-1.txt}\\[-0.5em]
\phantomsection\label{foro-participacion}

\textbf{Asunto:} La organización como sistema humano y abierto\par
\medskip
Hola, compañeras y compañeros. Comparto mi aportación sobre los rasgos que distinguen a una organización y los elementos necesarios para estudiarla.
\medskip

\begin{enumerate}[leftmargin=1.6em,labelsep=0.6em,itemsep=0.7em,label=\textbf{\arabic*)}]
  \item \textbf{¿Qué es una organización?} Una organización es un sistema social deliberadamente coordinado. Está integrado por personas y grupos que distribuyen responsabilidades, emplean recursos y realizan procesos para alcanzar propósitos comunes con cierta continuidad. No se limita al organigrama: también comprende relaciones informales, cultura, tecnología, decisiones y vínculos con el entorno. Como sistema abierto, recibe recursos e información, los transforma y genera resultados; la retroalimentación le permite aprender y adaptarse \citep{kastRosenzweig1972,kastRosenzweig1985}.

  \item \textbf{¿Qué no debe ser considerado una organización?} Una simple multitud, una fila o varias personas que coinciden temporalmente no constituyen una organización si carecen de propósito compartido, roles, coordinación y continuidad. Tampoco basta con tener un nombre, un edificio o un registro legal. En sentido contrario, no sería correcto afirmar que solo las empresas formales o lucrativas son organizaciones: una cooperativa, una asociación vecinal o un colectivo pueden serlo cuando articulan objetivos, reglas y acciones estables.

  \item \textbf{¿Cuáles son los principales elementos que debemos considerar al hacer estudios organizacionales?} Deben analizarse el propósito y los resultados esperados; las personas, los grupos, el liderazgo, la cultura y las relaciones de poder; la estructura, la división del trabajo y la coordinación; los procesos, la tecnología, la información y los recursos; y el ambiente económico, social, político y tecnológico. Estas dimensiones son interdependientes: una tecnología puede mejorar un proceso, pero fracasar si contradice la cultura, concentra decisiones o no desarrolla capacidades en el personal. También conviene observar entradas, transformación, resultados y retroalimentación, pues no existe una estructura ideal para todos los contextos \citep{bright2019,openstaxEnvironment2019,openstaxStructures2019}.
\end{enumerate}
\medskip
Mi postura es que estudiar una organización exige comprender cómo coordina diferencias y cómo aprende de sus efectos, no solamente describir puestos. Para dialogar: ¿qué elemento suele ignorarse más al diagnosticar una organización y qué consecuencia produce esa omisión?\par
\medskip
Saludos y quedo atento a sus comentarios.
\medskip
\textbf{Duración estimada del audio:} entre 2 minutos 35 segundos y 2 minutos 55 segundos.
\end{forobox}

\subsection{Interacción con dos aportaciones}

Las siguientes réplicas se adaptan a la idea concreta publicada por cada compañero; no afirman que la interacción ya ocurrió.

\begin{forobox}[title={Primera retroalimentación para adaptar y publicar}]
\textbf{Asunto:} Propósito y coordinación\par
\medskip
Coincido en que el objetivo común permite reconocer una organización y añadiría que dicho objetivo necesita traducirse en roles, decisiones y mecanismos de coordinación. Sin esa conexión, el propósito permanece como intención y no como acción colectiva. Desde el enfoque sistémico, también importa observar si los resultados obtenidos regresan como información para corregir el proceso \citep{kastRosenzweig1972}. ¿Qué mecanismo de coordinación consideras más decisivo en el ejemplo que planteas y por qué?
\end{forobox}

\begin{forobox}[title={Segunda retroalimentación para adaptar y publicar}]
\textbf{Asunto:} Ambiente y cultura\par
\medskip
Tu énfasis en las personas ayuda a evitar que la organización se reduzca a un organigrama. Complementaría la idea considerando cultura y ambiente: una misma estructura puede producir resultados distintos según las normas compartidas, el poder y las presiones externas. Por eso el diseño debe guardar congruencia con la estrategia y el contexto \citep{openstaxEnvironment2019,openstaxStructures2019}. ¿Qué señal permitiría detectar que la cultura está dificultando la adaptación de esa organización?
\end{forobox}

\subsection{Verificación previa a la publicación}

\begin{itemize}
  \item Ensayar el guion con cronómetro y mantener la grabación por debajo de tres minutos.
  \item Vocalizar las tres respuestas, sin leer las referencias bibliográficas completas.
  \item Publicar el audio mediante la herramienta disponible y comprobar que el enlace tenga acceso.
  \item Personalizar cada réplica con una idea real de la aportación correspondiente.
  \item Publicar al menos dos respuestas diferenciadas y conservar evidencia de la interacción.
\end{itemize}

\clearpage
\section{Conclusiones}

Una organización no se define por su edificio, lucro o formalización legal, sino por la coordinación relativamente estable de personas, recursos y procesos alrededor de propósitos. Su estudio requiere observar simultáneamente estructura, comportamiento humano, cultura, poder, tecnología y ambiente, porque una intervención en cualquiera de esos componentes modifica el sistema completo. La postura más útil para la gestión es, por tanto, relacional y contingente: diagnosticar conexiones, tensiones y retroalimentación antes de recomendar cambios. En el foro, esta delimitación permite pasar de definiciones generales a una conversación sobre consecuencias concretas y comparar perspectivas mediante las dos réplicas. La responsabilidad sobre la selección de fuentes, la argumentación y la publicación final permanece en el autor.\footnote{Declaración de uso: se utilizó inteligencia artificial como apoyo para organizar, revisar y editar el texto; el autor verificó las fuentes, asumió la postura expresada y conserva la responsabilidad sobre el contenido final.}

\clearpage
\bibliography{\itescabibliography}

\end{document}
